% A. Schrijver, Bounds on the number of Eulerian orientations,
% Combinatorica 3 (1983), 375--380, p. 379, Eq. (22).
% https://ir.cwi.nl/pub/10053/10053D.pdf
% The following is a transcription of the hypothesis and displayed conjecture.
% epsilon counts edge orientations with equal indegree and outdegree at every
% vertex; K_t is the complete undirected graph on t vertices.
if $G$ is a simple undirected graph with
degree sequence $2k_1,\ldots,2k_n$, then
\[
  \varepsilon(G)\leq
  \prod_{i=1}^n \varepsilon(K_{2k_i+1})^{1/(2k_i+1)}.
  \tag{22}
\]

% Srinjoy Srimani, Jacques Verstraëte and Jason Yang,
% Counting Eulerian Orientations in Regular Graphs, arXiv:2611.29726v1,
% section 1.2, Conjecture B (verbatim TeX), and section 7 (degree-six question).
% https://arxiv.org/abs/2611.29726v1
\begin{conjecture} {\bf (Schrijver~\cite{Sch83})} \label{conj:schrijver-degrees}
Let $G$ be an $n$-vertex graph with even degrees $d_1,d_2,\dots,d_n$, with no loops or multiple edges. Then
\[
\varepsilon(G)\le \prod_{i=1}^n \varepsilon(K_{d_i+1})^{1/(d_i+1)}.
\]
\end{conjecture}

For $d=6$ the
question is whether $\varepsilon(G)\le\varepsilon(K_7)^{n/7}$ for every
simple $6$-regular graph. Disjoint copies of $K_7$ attain this bound when
$7$ divides $n$.
